\section{Strong irreducibility and the automorphic type-I tangent space}
\label{sec:type-one-rigidity}

Let \(K/\mathbb Q\) be the imaginary quadratic field,
\(H=G_K\), and \(\epsilon=\epsilon_\ell\). In a normalized type-I
point \(r\), let \(\rho:H\to {\rm GL}_3(B)\) be its matrix part,
where \(B=\overline{\mathbb Q}_\ell\). The type-I condition is that
\(\rho\) has no proper nonzero invariant \(B\)-subspace. Write
\(\Psi=\epsilon^e\psi\), with \(e\ge1\), \(\psi\) of finite
cyclotomic conductor dividing \(N\). The ordinary and residual
conditions on the point supply an integral model with the following
two properties:
\begin{enumerate}
\item On decomposition at \(\ell\) inside \(H\) there is a common
invariant full flag; on inertia its successive eigencharacters
are \(\Psi,1,\Psi^{-1}\), in the order fixed by the window.
\item Reduction at the maximal ideal of the integral model satisfies
\[
 \overline{{\rm tr}\,\rho(s)}
   =\overline{\Psi(s)}+1+\overline{\Psi(s)}^{-1}
           \qquad(s\in H).
\]
\end{enumerate}
The residue hypothesis \(\ell-1\nmid 2e\) is essential to the
argument below. We first prove the open-subgroup irreducibility
assertion used in the original release.

\paragraph{Proposition.}
For every open subgroup \(V\) of \(H\), \(\rho|_V\) is irreducible.

\paragraph{Proof.}
Take an open normal subgroup \(U\) of \(G_\mathbb Q\) contained
in \(V\), obtained by intersecting its finitely many conjugates.
It is enough to prove irreducibility on \(U\). Choose a place over
\(\ell\). Its inertia subgroup lies in \(H\), since \(\ell\) does
not divide the discriminant of \(K\). The cyclotomic character on
this inertia subgroup surjects onto \(\mathbb Z_\ell^\times\):
the extensions \(\mathbb Q_\ell(\mu_{\ell^j})/\mathbb Q_\ell\)
are totally ramified with Galois groups
\((\mathbb Z/\ell^j\mathbb Z)^\times\), and passage to the inverse
limit gives the assertion. Choose \(\sigma\) on inertia with
\(\epsilon(\sigma)=1+\ell\). Put
\(m=[G_\mathbb Q:U]\cdot|(\mathbb Z/N)^\times|\). Then
\(u=\sigma^m\) lies in \(U\), and \(\psi(u)=1\). The three
eigenvalues of \(\rho(u)\) are distinct since
\(\Psi(u)^2=(1+\ell)^{2em}\ne1\).

Choose also \(x\) on inertia whose cyclotomic character generates
\(\mathbb F_\ell^\times\) after reduction. The window assumes
\(\psi^2=1\) on this inertia subgroup. Consequently
\(\overline{\Psi(x)}^2=\overline{\epsilon(x)}^{2e}\ne1\).
Both \(x\) and \(u\) preserve the ordinary flag. Each subspace
in the flag is spanned by some eigenlines of the regular semisimple
element \(\rho(u)\): its spectral projections are polynomials in
\(\rho(u)\).

Suppose \(\rho|_U\) is reducible. The restriction is semisimple:
take a nonzero simple \(U\)-submodule and all its \(H\)-translates.
Normality of \(U\) makes each translate simple, and their sum is
the entire space by \(H\)-irreducibility. A sum of simple
submodules is semisimple. The group \(H\) permutes its isotypic
components transitively, so they have equal dimension and equal
multiplicity of their respective simple constituents.
As the space has dimension three, reducibility on \(U\) now gives
two possibilities. One is three copies of one character, on which
\(u\) would be scalar. This contradicts its distinct eigenvalues.
The other is three pairwise distinct characters, with three lines
permuted transitively by \(H\). These lines are the eigenlines of
\(\rho(u)\). The permutation induced by \(x\) preserves each stage
of the complete ordinary flag and hence fixes the three lines
individually.

The transitive permutation image of \(H\) on three letters has
order divisible by three. Choose \(g\in H\) inducing a three-cycle.
Then \(\rho(gx^j)\) has zero diagonal in the eigenbasis of
\(\rho(u)\), for every \(j\ge0\). Its trace therefore vanishes.
In the residue field put \(a=\overline{\Psi(g)}\) and
\(b=\overline{\Psi(x)}\). The residual trace formula says
\[
 a b^j+1+(a b^j)^{-1}=0\qquad(j\ge0).
\]
Thus all the nonzero elements \(a b^j\) are roots of
\(T^2+T+1\). In particular among \(a,ab,ab^2\) two coincide,
forcing \(b=1\) or \(b^2=1\). This contradicts the choice of \(x\).
Irreducibility on \(U\) and on \(V\) follows. \(\square\)

\paragraph{Lemma.}
Let \(V\) be a finite-dimensional characteristic-zero continuous
representation of \(H\) over \(B\) which is irreducible on every
open subgroup of \(H\). For any continuous abelian quotient
\(H\to C\), the restriction of \(V\) to its kernel is irreducible.

\paragraph{Proof.}
Let \(M\) be the Zariski closure of the image in \({\rm GL}(V)\);
write \(M^0\) for its identity component. The preimage of \(M^0\)
is open in \(H\): there are finitely many components, each closed
in the \(\ell\)-adic topology, so its complement is closed as well.
Its image is Zariski dense in \(M^0\), since this component is also
open in \(M\) for the Zariski topology. In particular \(M^0\)
acts irreducibly. Its unipotent radical has a nonzero fixed vector
by triangularization. The whole fixed space is \(M^0\)-stable by
normality, so the unipotent radical fixes \(V\) pointwise and is
trivial in this faithful realization. It follows that \(M^0\) is
reductive. Its connected center acts by scalars by Schur's lemma.
A connected reductive group is the product of its connected center
and its derived subgroup. Hence \((M^0)^{\rm der}\) already acts
irreducibly on \(V\).

All group commutators of the preimage of \(M^0\) belong to the
kernel of \(H\to C\). The Zariski closure of its image there
contains the commutators of all pairs in \(M^0\), by density
and the algebraicity of the commutator map. It therefore
contains \((M^0)^{\rm der}\). A subspace invariant under the kernel
is invariant under its Zariski closure, and the preceding
irreducibility finishes the proof. \(\square\)

\paragraph{The polarized adjoint.}
In addition assume the type-I automorphy conclusion: \(\rho(-1)\)
is attached to a regular algebraic, conjugate self-dual cuspidal
automorphic representation \(\pi\) of \({\rm GL}_3(\mathbb A_K)\).
This is the automorphic witness in \(\mathrm{IsAutIAlt}\), whose
definition also carries local compatibility information. Attachment
compares the unramified normalized Satake roots with geometric
Frobenius; its factor for rank three is \(({\rm N}v)^{(3-1)/2}\).
With this normalization the semisimple representation
\(r_{\pi,\iota}\) of the automorphic Galois theorem identifies with
\(\rho(-1)\) by Chebotarev and Brauer--Nesbitt. The preceding lemma
for the cyclotomic quotient shows that \(\rho\), and hence its scalar
twist, is irreducible on \(G_{K(\mu_{\ell^\infty})}\).

Write \(J=J_3\) for the antidiagonal identity matrix, and choose
the normalized complex conjugation \(c\), whose matrix part is \(1\).
The semidirect-product law of the point gives
\[
 \rho(chc^{-1})=J\,{}^t\!\rho(h)^{-1}J^{-1}\qquad(h\in H).
\]
For \(V=\rho(-1)\), \(\epsilon(chc^{-1})=\epsilon(h)\) gives
\[
 V(chc^{-1})=\epsilon(h)^{-2}J\,{}^t\!V(h)^{-1}J^{-1}.
\]
Thus \(J\) supplies the symmetric polarization
\(V^c\simeq V^\vee\otimes\epsilon^{1-3}\).
Any two such polarizations on the irreducible representation
differ by a scalar, which cancels in the adjoint. On \(H\) the
polarized adjoint acts by conjugation by \(V(h)\), or equivalently
by \(\rho(h)\); on \(c\) it acts on all nine matrix entries by
\(X\mapsto-J\,{}^t\!XJ^{-1}\). These are precisely the actions
in the tangent module of the point.

In \emph{On the vanishing of adjoint Bloch--Kato Selmer groups of
irreducible automorphic Galois representations}, Theorem 6.1,
J.\ A.\ Thorne proves that for a CM field \(F\), a polarizable
regular algebraic cuspidal \(\pi\) over \(F\), and
\(r_{\pi,\iota}\) irreducible over \(F(\mu_{\ell^\infty})\), the
Selmer group of the polarized adjoint over \(F^+\) vanishes.
With \(F=K\) the base \(F^+\) is \(\mathbb Q\). The local
conditions are unramified away from \(\ell\) and the crystalline
kernel at \(\ell\), as in Section~\ref{sec:kato-rigidity}.
We have checked the polarization and the irreducibility, so
\[
 H^1_f(\mathbb Q,{\rm ad}\,r)=0.
\]
This applies after the type-I automorphic witness has been obtained.

The original route in \(\mathrm{rigIUE\_alt}\) factors the local
point through an integral valuation ring, uses the ordinary local
conditions, establishes the open-subgroup assertion
\(\mathrm{matIrr\_open\_gd}\), and obtains a free-ratio condition
from the automorphic twist. The subsequent result
\(\mathrm{rigEnd3\_alt}\) passes to the auxiliary CM field of the
frame, transports automorphy and irreducibility there, applies
adjoint lifting, and descends the vanishing. The propositions above
give the open-subgroup argument and a direct application of the
published theorem on the original quadratic field. The existence of
the required automorphic representation is an input to both versions.
